\chapter{Preface}

\jns{} is a framework for automated reaction-path search.  It drives an
open-source reaction-network engine (SCINE Chemoton and Puffin, from the
Reiher group) with machine-learning potentials (MACE and AIMNet2), and it
keeps track of how much the potentials can be trusted: a committee of models
from different architecture families measures its own disagreement on every
geometry, and predictions beyond a calibrated uncertainty threshold are
handed back to a higher-level quantum-chemistry program (ORCA).

This manual documents version 0.1.0.  It is written for two kinds of
readers: users who want to run explorations with \jns{} (Chapters
\ref{ch:installation}--\ref{ch:configuration}), and developers who want to
extend or integrate it (Chapters \ref{ch:architecture} onward, and the
appendices).

\section*{How to read this manual}

\begin{readbox}[How to read this manual]
\begin{itemize}
  \item If you have never run \jns{}: start with Chapter
    \ref{ch:installation} (installation) and Chapter
    \ref{ch:getting-started} (first five minutes).  The calibration example
    in Chapter \ref{ch:getting-started} runs with nothing but NumPy.
  \item If you want to put the NNP into a running exploration: Chapter
    \ref{ch:scine-integration} describes the bridge contract and the
    injection patches, including the exact deployment steps.
  \item If you want to understand (or audit) the uncertainty machinery:
    Chapter \ref{ch:uncertainty} defines the disagreement measures and the
    calibration protocol in full.
  \item Reference material -- calculators, scripts, configuration fields,
    constants, environment variables -- lives in Chapters
    \ref{ch:calculators}, \ref{ch:scripts}, \ref{ch:configuration} and
    Appendices \ref{app:constants}--\ref{app:glossary}.
\end{itemize}
\end{readbox}

\section*{Conventions}

\begin{itemize}
  \item Energies and forces follow the ASE convention throughout: energies
    in electronvolts (eV), forces in eV/\AA{}; reaction barriers are quoted
    in kcal/mol where they come from the benchmark protocols.
  \item Atomic units (Hartree, Bohr) appear only where a program's native
    output uses them; the conversion constants are collected in Appendix
    \ref{app:constants}.
  \item \code{monospace} marks code, file names and identifiers;
    \texttt{transcript} blocks show terminal sessions and
    \texttt{codeblock} blocks show code.
  \item The product name is written \jns{}; the Python package is
    \code{jnuscout}, the repository is
    \code{github.com/jqChen1566/jnu-scout}.
\end{itemize}

\section*{Authors, funding and citation}

\begin{readbox}[Authors and funding]
\jns{} is developed by Jianqi Chen, Zhuoran Zhang and Laiyu Zhang (College
of Chemistry and Materials Science, Jinan University, Guangzhou, China),
Mingyu Chen (Information Hub, The Hong Kong University of Science and
Technology (Guangzhou), China) and Ruiqian Luo (School of Chemistry,
Sun Yat-sen University, Guangzhou, China), under the supervision of
Prof.\ Juan Li (Jinan University).

The project is supported by the Jinan University Provincial College
Students' Innovation and Entrepreneurship Training Program
(Project No.\ S202610559049).
\end{readbox}

If you use \jns{} in academic work, please cite it as described in the
\file{CITATION.cff} file of the repository; the same information appears in
the README.

\section*{License}

\jns{} is distributed under the Apache License 2.0 (see \file{LICENSE} and
\file{NOTICE}).  It builds on external open-source and academic software --
SCINE, MACE, AIMNet2, ASE, the ORCA program and the xtb executable -- none
of which is bundled with this distribution; each remains under its own
license and must be installed separately.
